\section{Devroye's series, the mean excess function and detection}\label{ev:sec:devroye}

This section tests the extreme value picture through the Laplace transform $L(s)=\xi(\frac12)/\xi(\frac12+\sqrt s)$ and a little information about the zeros: Devroye's alternating series, the mean excess function with Schoenberg's criterion, the weak convergence behind them, and detection in the de Bruijn Newman family. Cited, not restated: the residues $R_k=c_k\gamma_k^2$ and the spectral expansion of the density (Section~\ref{sec:ch18:law} and Exercise~\ref{ex:ch18:residues}), the criterion by P\'olya frequency functions (Proposition~\ref{tx2:prop:pf}), the weak closure of the GGC class (Corollary~\ref{cor:weak:closure}) and the digits rule for Prony fits (Section~\ref{sec:ch18:reach}).

\subsection{Devroye's alternating series and the sign pattern}\label{ev:sec:sign}

Devroye's exact samplers for the Jacobi theta laws rest on one device: a density written as a series whose terms alternate in sign and decrease in size is bracketed by its partial sums \cite{Devroye09,Devroye86}. % UNCHECKED: Devroye 2009 (Statist. Probab. Lett. 79), the specific content of the paper was not read; the alternating series method itself is in Devroye86.
The reciprocal $\Xi$ law has such a series under hypotheses. Assume RH, simple zeros and $|\xi'(\rho)|\ge e^{-C|\gamma|}$ as in Theorem~\ref{ht:thm:res}; the zeros are $\rho=\frac12+i\gamma_k$ with $\gamma_k$ real, $0<\gamma_1<\gamma_2<\cdots$, and $\pm\gamma_k$ are the zeros of $\Xi$. Then the density of $X=\sum_kE_k/\gamma_k^2$ is
\[
f(x)=\sum_{k\ge1}c_k\gamma_k^2e^{-\gamma_k^2x},\qquad c_k=\prod_{j\ne k}\Big(1-\frac{\gamma_k^2}{\gamma_j^2}\Big)^{-1}=\frac{2i\,\xi(\frac12)}{\gamma_k\,\xi'(\frac12+i\gamma_k)} .
\]
Without the lower bound for $|\xi'(\rho)|$ only the grouped limit of the residue sums is proved (Theorem~\ref{ht:thm:grouped}), and the partial sums below are not known to converge to $f$; at a multiple zero $c_k$ is undefined. The two term tail of $f$ needs none of these hypotheses (Proposition~\ref{ht:prop:tail}).

\begin{proposition}[Sign pattern]\label{ev:prop:sign}
If the atoms $\gamma_1^2<\gamma_2^2<\cdots$ are real and simple, then $\operatorname{sign}c_k=(-1)^{k-1}$.
\end{proposition}

\begin{proof}
Exactly $k-1$ of the factors $1-\gamma_k^2/\gamma_j^2$ are negative, those with $j<k$.
\end{proof}

\begin{remark}
If a pair of atoms $\lambda_k=\gamma_k^2$ is not real, the two coefficients of the pair are complex conjugates, and for every real atom $\lambda_m$ the pair contributes the factor $|1-\lambda_m/\lambda_k|^{-2}>0$ to $c_m$, so the coefficient of a real atom of rank $r$ among the real atoms has sign $(-1)^{r-1}$. A sign test therefore sees a nonreal pair only if the rank of each atom is known, which is a zero count. The test also reads $\gamma_k$ and $\xi'(\rho_k)$, not $L$ alone.
\end{remark}

Devroye's switch point is the $x_k$ beyond which $|c_{k+1}|\gamma_{k+1}^2e^{-\gamma_{k+1}^2x}\le|c_k|\gamma_k^2e^{-\gamma_k^2x}$, namely
\[
x_k=\frac{\log\big(|c_{k+1}|\gamma_{k+1}^2/(|c_k|\gamma_k^2)\big)}{\gamma_{k+1}^2-\gamma_k^2}.
\]
Bracketing at a point $x$ needs the terms to alternate and decrease from the first index on, that is $x\ge\sup_kx_k$.

\begin{table}[h]\centering\small
\begin{tabular}{lll}
\toprule
law & signs of $c_k$ & $x_1,x_2,x_3$\\
\midrule
$1/\cosh\sqrt{2s}$ & $+-+-+-+-$ & $.111,\ .0259,\ .0114$\\
$\sqrt{2s}/\sinh\sqrt{2s}$ & $+-+-+-+-$ & $.0936,\ .0329,\ .0167$\\
reciprocal $\Xi$ & alternating, first $1700$ zeros & $.01963,\ .01533,\ .01385$\\
\bottomrule
\end{tabular}
\caption{Calibration on Devroye's two explicit laws, with atoms $\pi^2(k-\frac12)^2/2$ and $\pi^2k^2/2$, and the reciprocal $\Xi$ law with atoms $\gamma_k^2$. For the two strings $c_k\lambda_k=(-1)^{k-1}\pi(k-\frac12)$ and $c_k=2(-1)^{k-1}$ in closed form, so that $x_k=\log\frac{2k+1}{2k-1}/(k\pi^2)$ and $x_k=\log\frac{(k+1)^2}{k^2}/\big((2k+1)\pi^2/2\big)$ (computed, and checked against the product at $K=4000$ factors). For the reciprocal $\Xi$ law the coefficients come from $\xi'(\rho_k)$ at the first $1700$ zeros (computed).}\label{ev:tab:devroye}
\end{table}

For the reciprocal $\Xi$ law the signs alternate for all $1700$ coefficients (computed): $c_1=50.871$, $c_2=-2664.6$, $c_3=3.1420\times10^4$, $c_4=-1.3550\times10^6$, and so on. Beyond the first switch point the partial sums bracket the density, checked against Talbot inversion at 30 digits.
\begin{itemize}
\item At $x=0.02$, where $f=80.10420$, the partial sums are $186.93$, $16.148$, $88.602$, $77.170$, $80.2026$, $80.09549$.
\item At $x=0.03$, where $f=23.433069$, they are $25.352$, $23.295$, $23.4341$, $23.43301$, $23.4330693$, $23.43306926$.
\end{itemize}
Each sum lies on the opposite side of $f$ from the one before. Over the first $1700$ zeros the largest switch point is $x_1=0.019627$, and the largest from $k=4$ on is $0.011656$, at $k=7$ (computed). This gives an exact Devroye sampler for $X$ on $x\ge0.01963$, and two sided bounds there, conditional on RH with simple zeros, on the lower bound for $|\xi'(\rho)|$, and on $x_k\le0.01963$ for every $k$, which was checked for $k\le1699$ only. For large $k$ the switch point depends on the ratios $|\xi'(\rho_k)/\xi'(\rho_{k+1})|$ over the gaps $\gamma_{k+1}^2-\gamma_k^2$, for which no bound is known. The bounds are computed in floating point, not enclosed.

\subsection{Schoenberg's theorem and the mean excess function}\label{ev:sec:mexcess}

The mean excess (mean residual life) function of $X$ is $e(u)=\E[X-u\mid X>u]$. It can be computed without zeros, by one Laplace inversion each for numerator and denominator,
\[
\int_u^\infty\mathbb{P}(X>x)\,dx=\mathcal L^{-1}\Big[\frac{\E X}{s}-\frac{1-L(s)}{s^2}\Big](u),
\]
with
\[
\E X=\frac{\xi''(\frac12)}{2\xi(\frac12)}=0.02310499311541897\ldots
\]
The criterion by P\'olya frequency functions (Proposition~\ref{tx2:prop:pf}) rests on Schoenberg's characterisation \cite{Schoenberg51}, restated with its nondegeneracy condition and the Gaussian parameter renamed to avoid a clash with $\gamma_k$.

\begin{theorem}[Schoenberg]\label{ev:thm:schoenberg}
An integrable density $f$ is a P\'olya frequency function, totally positive of every order, exactly when its bilateral Laplace transform is $1/\Psi(s)$ in a strip around the origin, with $\Psi(s)=Ce^{-\alpha s^2+\delta s}\prod_k(1+\delta_ks)e^{-\delta_ks}$, where $\alpha\ge0$, $\delta$ and $\delta_k$ are real, $\sum\delta_k^2<\infty$ and $0<\alpha+\sum\delta_k^2$ (Schoenberg 1951).
\end{theorem}

\begin{corollary}\label{ev:cor:pf}
RH holds if and only if $L(s)=\xi(\frac12)/\xi(\frac12+\sqrt s)$ is the Laplace transform of a P\'olya frequency density. Under RH that density is log concave, the hazard rate of $X$ is increasing, and the mean excess function decreases monotonically,
\[
e(0)=\sum_k\gamma_k^{-2}=0.0231050\ \searrow\ e(\infty)=\gamma_1^{-2}=0.00500524 .
\]
\end{corollary}

\begin{proof}
The function $1/L(s)=\prod_k(1+s/\gamma_k^2)$ has the form of Theorem~\ref{ev:thm:schoenberg} with $\alpha=0$ and $\delta_k=\gamma_k^{-2}$, which are real exactly when every $\gamma_k$ is real: an off line zero with $\gamma_k$ not real gives a nonreal $\gamma_k^2$, and a real zero of $\xi$ off the line would lie in $(0,1)$, where, by the classical fact that $\zeta$ has no zero in $(0,1)$, there is none. Total positivity of order two is log concavity, which implies an increasing hazard rate, and an increasing hazard rate implies a decreasing mean residual life (Barlow and Proschan \cite{BarlowProschan75}). % UNCHECKED: Barlow and Proschan, the implications are standard, the book was not read.
\end{proof}

So the mean excess function interpolates between the Thorin moment $\sum_k\gamma_k^{-2}$ at $u=0$ and the first atom at $u=\infty$, and its approach to the limit carries the next atom; the analogous hazard limit, with rate $\gamma_2^2-\gamma_1^2$, is in Section~\ref{sec:ch18:law} and Exercise~\ref{ex:ch18:hazard}.

\begin{proposition}[Expansion and rate]\label{ev:prop:meexp}
Under RH with simple zeros and $|\xi'(\rho)|\ge e^{-C|\gamma|}$ as above, and with $\lambda_k=\gamma_k^2$,
\[
\begin{aligned}e(u)&=\frac{\sum_kc_ke^{-\lambda_ku}/\lambda_k}{\sum_kc_ke^{-\lambda_ku}}=\frac1{\lambda_1}+A\,e^{-(\lambda_2-\lambda_1)u}\\&\quad+O\big(e^{-\min(\lambda_3-\lambda_1,\,2(\lambda_2-\lambda_1))u}\big),\qquad A=\frac{c_2}{c_1}\Big(\frac1{\lambda_2}-\frac1{\lambda_1}\Big)>0 .\end{aligned}
\]
The excess law of $X$ over $u$ converges to the exponential law with rate $\lambda_1$ in total variation at rate $e^{-(\lambda_2-\lambda_1)u}$. Without simple zeros the two term form follows from the two term tail of Proposition~\ref{ht:prop:tail}, with $w_k/\lambda_k$ in place of $c_k$ for $k=1,2$.
\end{proposition}

\begin{proof}
The first equality is $\int_u^\infty\mathbb{P}(X>x)\,dx\big/\mathbb{P}(X>u)$ with $\mathbb{P}(X>x)=\sum_kc_ke^{-\lambda_kx}$. Put $\varepsilon=e^{-(\lambda_2-\lambda_1)u}$. The numerator is $\lambda_1^{-1}c_1e^{-\lambda_1u}\big(1+\frac{c_2\lambda_1}{c_1\lambda_2}\varepsilon+O(e^{-(\lambda_3-\lambda_1)u})\big)$ and the denominator is $c_1e^{-\lambda_1u}\big(1+\frac{c_2}{c_1}\varepsilon+O(e^{-(\lambda_3-\lambda_1)u})\big)$. Dividing, the first order term is $\frac{c_2}{c_1}\big(\frac1{\lambda_2}-\frac1{\lambda_1}\big)\varepsilon$ and the next is $O(\varepsilon^2)$. Since $c_2/c_1<0$ (Proposition~\ref{ev:prop:sign}) and $1/\lambda_2<1/\lambda_1$, $A>0$.
\end{proof}

The following were computed from the residue series over the first $300$ zeros at 50 and 70 digits on the grid $u=0.005,0.0075,\ldots,0.2$ (79 points); the source used Talbot inversion from $L$ alone. Nothing is enclosed in interval arithmetic.
\begin{itemize}
\item Log concavity: $ff''-f'^2<0$ at all 79 points (a sampled check). Two of the points, $0.005$ and $0.0075$, lie below $0.008$, where Theorem~\ref{ht:thm:lc} starts and where the book's proof is conditional on a saddle point approximation (Proposition~\ref{ht:prop:lc0}).
\item Monotonicity: $e(u)$ decreases strictly at every grid step. Beyond $u\approx0.165$ consecutive values agree with $\gamma_1^{-2}=0.00500524412359412$ to every printed digit.
\item Expansion: Proposition~\ref{ev:prop:meexp} predicts $A=0.1436441604$. The observed deviation $e(u)-\gamma_1^{-2}$ against $Ae^{-(\lambda_2-\lambda_1)u}$ is $9.808\times10^{-5}$ against $1.006\times10^{-4}$ at $u=0.03$, $2.3625\times10^{-7}$ against $2.3637\times10^{-7}$ at $u=0.055$, and $1.3052007\times10^{-12}$ against $1.3052008\times10^{-12}$ at $u=0.105$. From $u=0.13$ on the agreement is to all eight printed digits; the residual at small $u$ is the $\lambda_3$ term.
\end{itemize}
Read without zeros, as the source does, the expansion gives the first two zeros. An Aitken extrapolation of three values of $e(u)$ estimates $e(\infty)=\gamma_1^{-2}$ and, as $-\log r/h$ for the extrapolated ratio $r$ and the step $h=0.0075$, the gap $\lambda_2-\lambda_1$. It gives $\gamma_1=14.1347778$, $14.1347266$ and $14.1347252$ from windows starting at $u=0.03$, $0.04$ and $0.05$, and $\gamma_2=20.965$, $21.006$ and $21.019$, against $\gamma_1=14.1347251417$ and $\gamma_2=21.0220396$. This is weaker than the Prony fits to the tail (Table~\ref{ev:tab:prony2}); the mean excess is a criterion and a check, not a way of computing zeros.

\subsection{Weak convergence}\label{ev:sec:weak}

Three weak convergence statements sit behind these diagnostics.
\begin{enumerate}
\item Pickands, Balkema and de Haan: the excess law over $u$, rescaled by $e(u)$, converges to a generalised Pareto law \cite{Pickands75,BalkemaDeHaan74}. % UNCHECKED: Pickands 1975 and Balkema and de Haan 1974, the standard statement is quoted, the papers were not read.
Here the shape is zero and the convergence is exponentially fast (Proposition~\ref{ev:prop:meexp}).
\item Closure: generalised gamma convolutions are closed under weak limits (Bondesson \cite{Bondesson}; Corollary~\ref{cor:weak:closure}), and so are P\'olya frequency laws with nondegenerate limits (Karlin \cite{Karlin68}). % UNCHECKED: Karlin, Total Positivity, the specific statement on nondegenerate limits was not read.
Hence, by Proposition~\ref{ht:prop:red} and Corollary~\ref{cor:weak:closure}, RH is equivalent to the reciprocal $\Xi$ law being a weak limit of hypoexponential laws $\sum_{k\le K}E_k/\lambda_{k,K}$ with real positive rates, a distributional form of the Hurwitz argument.
\item The exponential B splines of Proposition~\ref{ht:prop:bspline} (Section~\ref{ev:sec:bspline}) are such a sequence, but they use the zeros. Zero free sequences exist under RH: the Jensen laws of Section~\ref{ht:sec:jensen} are GGC laws for every degree $d\le2\cdot10^{17}$, exist for every $d$ exactly when RH holds, and then converge to the law of $X$ (Proposition~\ref{ht:prop:jensen}). Another candidate is a sequence of approximations $\widehat L_N(s)=\widehat H_N(0)/\widehat H_N(2i\sqrt s)$ with only real zeros that converge to $L$. P\'olya's kernels give P\'olya frequency laws but do not converge. The folded kernels converge, and their zeros were found to be real up to a height of about $c(N+1)^2$ with computed constants $c$ between $3.6$ and $4.4$ (certified in ball arithmetic for $N\le4$ given the count of zeros of $\Xi$ in each disc, computed for $N=5$; Sections~\ref{ev:sec:work} and~\ref{ev:sec:round2}).
\end{enumerate}

\subsection{Detection in the de Bruijn and Newman family}\label{ev:sec:detect}

The family $H_t$ is that of Chapter~\ref{ch:dbn}, with the variable of the zeros $w=z/2$; its collisions and Hankel test are in Sections~\ref{sec:ch21:collisions} and~\ref{sec:ch21:hankel}. At $t=-4$ the fourth and fifth zeros have merged, according to the source (not checked here). By quadrature against the kernel, $H_{-4}(2w)$ is negative at all 41 points of $[28,38]$ with step $0.25$ (computed; the largest value is $-6.67\times10^{-12}$ and the smallest modulus is at $w=33.25$), so it almost touches zero near $33$. Root finding locates the pair at
\[
w=33.1150940\pm0.1505390\,i,\qquad\gamma^2=1096.58679\pm9.970228\,i
\]
(computed). A short window diagnostic should miss it: the pair adds $e^{-1096.6x}\cos(9.97x)$ to the tail, and over a window of length $0.04$ the cosine turns by only $0.4$ radians. The source reports, and this was not recomputed, that a six atom Prony fit on such a window returns spurious atoms, one real and a pair, and that the same fit at $t=0$ returns an equally spurious pair where the true zero is real. % UNCHECKED: the spurious atoms of the short window fits (32.28, 34.78 +- 1.71 i, 33.32 +- 2.33 i) were not recomputed; the values are not printed.
Short windows cannot separate a slightly nonreal pair from omitted atom bias.

\begin{remark}[Detection rule, heuristic]\label{ev:rem:detect}
To see one radian of oscillation from a nonreal atom $\lambda=a+ib$, a tail window must extend to $x\approx1/b$. There its term is smaller than the leading term by $e^{-(a-\lambda_1)/b}$, so that the computation needs about $(a-\lambda_1)/(b\ln10)$ digits beyond the size of the leading term. If the window ends at $x_{\mathrm{end}}$ instead, about $(a-\lambda_1)x_{\mathrm{end}}/\ln10$ digits are needed. Section~\ref{sec:ch18:reach} states the same kind of rule for real atoms. This is an order of magnitude estimate, not a theorem.
\end{remark}

At $t=-4$ the first count is $(1096.59-234.42)/(9.970\ln10)\approx37.6$ digits (computed, with $\gamma_1=15.31093766$ taken from Table~\ref{ev:tab:detect}, which is unchecked). The source's window, below, ends at $0.188$, where the pair is smaller than the leading term by about $e^{-162}\approx10^{-70}$, the count for that window.

The source tested the rule with a window $x_j=0.08+0.012j$, $j=0,\ldots,9$, which turns the cosine by about one radian, and five atom Prony fits at 120 digits, with $H_{-4}(2i\sqrt s)$ computed as the Gaussian average $\E H(2i\sqrt s+W)$, $W$ normal with variance $8$, by 90 point Gauss and Hermite quadrature on $\xi$. Table~\ref{ev:tab:detect} is the outcome as the source reports it; the fits were not run for this book, and the only entries recomputed are the pair, the zero of $H_{-4}$ found directly, and the control column, which is $\gamma_1,\ldots,\gamma_5$.

\begin{table}[h]\centering\small
\begin{tabular}{ccc}
\toprule
 & $t=0$ (control) & $t=-4$\\
\midrule
$\widehat\gamma_1$ & 14.13472514 & 15.31093766\\
$\widehat\gamma_2$ & 21.02203964 & 22.5998361\\
$\widehat\gamma_3$ & 25.01085758 & 26.14438156\\
$\widehat\gamma_4$ & 30.42487613 & $33.11509402-0.1505390225\,i$\\
$\widehat\gamma_5$ & 32.93506159 & $33.11509402+0.1505390225\,i$\\
\bottomrule
\end{tabular}
\caption{Long window threshold fits from the Laplace transform alone, as reported by the source (unchecked: the fits were not recomputed; the pair agrees with the zero of $H_{-4}$ found directly, $33.1150940172\pm0.150539022\,i$).}\label{ev:tab:detect}
\end{table}
% UNCHECKED: Table tab:detect, the long window Prony fits at lambda = 0 and -4 (PENDING in checks-G2, needs ev_detect.py); only the pair (zero of H_{-4}) and the control zeros were recomputed.

If the table is right, the fit returns $\gamma_1,\ldots,\gamma_5$ for $t=0$ and for $t=-4$ a conjugate pair with Thorin atom $1096.58679\pm9.970227764\,i$, agreeing with the root found directly to all ten digits printed. A threshold analysis of the tail of the inverse Laplace transform of $1/H_{-4}(2i\sqrt s)$ (at $t=-4$ this need not be a probability density, so it is not a hitting time tail) would then detect a genuine loss of reality from the transform alone, once the window and the precision follow Remark~\ref{ev:rem:detect}. Until the fits are recomputed this is a claim of the source, and the rule itself stays heuristic.

The rule quantifies the limitation of all tail diagnostics. A pair that has only just left the line has small $\operatorname{Im}\gamma^2=2\operatorname{Re}\gamma\operatorname{Im}\gamma$, and the digits required grow like $1/\operatorname{Im}\gamma^2$, so it is the hardest kind to see. This is the Lehmer pair phenomenon in tail form (the flow of a Lehmer pair is in Section~\ref{sec:ch21:bounds}), and it is why such tests are verification to a finite height and precision rather than proof.

\section{Remaining open problems}\label{ev:sec:open}
\begin{enumerate}
\item Interval arithmetic for the Rouch\'e test at $N=5$ and for the count of zeros of $\Xi$ in each disc, and the level $N=6$ near height $196$ (Sections~\ref{ev:sec:work} and~\ref{ev:sec:round2}).
\item Identify the continuum extension for $\mu=\tfrac54$ and the interlacing string for $\mu=\tfrac94$, and combine them into a string for P\'olya's two term kernel.
\item Give the Hankel order needed to see the de Bruijn and Newman collision at $t\approx-0.881$ near height $95.5$. This is answered in part by Table~\ref{tab:ch21:nfail}: no order up to $50$ fails on $[-1.03,0]$, so at the collision time the order needed exceeds $50$.
\item Prove that the off line zeros of the folded kernels lie on a curve, with explicit spacing, and decide whether masking persists at all Hankel orders for this family.
\item Turn the Prony scans into an error formula in $(K,x_0,\text{precision})$ and apply it to a Lehmer pair.
\item Find a product structure in the theta superposition of Gumbel pieces that carries the Thorin atoms at $\gamma_k^2$ (Section~\ref{ev:sec:bspline}).
\item Show that the switch points satisfy $x_k\le x_1$ for every $k$ (Section~\ref{ev:sec:sign}), from properties of $L$ and a lower bound for $|\xi'(\rho)|$, and use Devroye's bounds for a certified sampler of the reciprocal $\Xi$ law. The alternation itself is proved (Proposition~\ref{ev:prop:sign}); what is open is the monotonicity of the terms.
\item Higher order total positivity: test $\mathrm{PF}_3$ and beyond for the reciprocal $\Xi$ density from Laplace inversion, since Corollary~\ref{ev:cor:pf} makes every order a necessary condition for RH. The ladder of Toeplitz minors of the P\'olya density is in Chapter~\ref{chap:hitting}.
\item Make the detection rule of Remark~\ref{ev:rem:detect} into an algorithm with error control, and estimate its cost for the Lehmer pairs known from verification.
\item \begin{resultpar}\label{ev:res:dyckside}One sided approach in the Dyck chain. If $F_n\to2\xi$ locally uniformly on $\operatorname{Re}s>0$ and, for every $T>0$, the Dirichlet polynomial $F_n$ of the Dyck chain (Section~\ref{ev:sec:decomp}) has no zero in $\tfrac12<\operatorname{Re}s<1$, $0<\operatorname{Im}s<T$ for all large $n$, then RH holds: if RH fails, the functional equation and the symmetry $\xi(\bar s)=\overline{\xi(s)}$ give a zero $\rho$ of $\xi$ with $\tfrac12<\operatorname{Re}\rho<1$ and $\operatorname{Im}\rho>0$, and by Hurwitz's theorem, for $T>\operatorname{Im}\rho$, $F_n$ has zeros in that rectangle for all large $n$. Under Conjecture~\ref{ev:conj:lbias} the hypothesis holds for $T<\gamma_{64}=169.9119\ldots$ and fails for every $T>\gamma_{64}$. The zeros of $\xi$ below height $170$ are $\rho_1,\dots,\rho_{64}$, simple and on the line \cite{LRW86}. For $T>\gamma_{64}$, the zero $z_n(\rho_{64})$ of Proposition~\ref{ev:prop:lbias} tends to $\rho_{64}$ and has $\operatorname{Re}z_n(\rho_{64})=\tfrac12+\delta_n\operatorname{Re}c_{\rm D}(\rho_{64})+O(1/n)>\tfrac12$ for large $n$, since $\operatorname{Re}c_{\rm D}(\rho_{64})>53.08$ (computed in ball arithmetic, Section~\ref{ev:sec:decomp}). For $T<\gamma_{64}$, $F_n\to2\xi$ uniformly on the closed rectangle; outside small discs about $\rho_1,\dots,\rho_{63}$ it carries no zero of $\xi$, hence none of $F_n$ for large $n$, and in the disc about $\rho_k$ the only zero is $z_n(\rho_k)$, to the left of the line for large $n$ because $\operatorname{Re}c_{\rm D}(\rho_k)<0$ (computed in ball arithmetic). This is the argument of Corollary~\ref{rz:cor:Tfalse} for the Dyck chain: conditionally on Conjecture~\ref{ev:conj:lbias} the hypothesis fails, and one sided approach cannot prove RH along this chain. For $10^5\le n\le10^8$ the zero near $\rho_{64}$ lies to the right of the line (computed, Section~\ref{ev:sec:decomp}). Status: the implication proved, the failure of the hypothesis conditional on Conjecture~\ref{ev:conj:lbias}; open: prove the conjecture, which would also give the locally uniform convergence of the first hypothesis. This item is based on \cite{P26excursion,P26lbias}.\end{resultpar}
\end{enumerate}
